\section*{Abstract}

The measurement-based quantum computing (MBQC) model presents a way to do quantum computations in constant depth by measuring a pre-prepared graph state. As such, there is interest in determining which graph states are universal for MBQC. The universality of the comparability grid is currently an open problem, and of some interest as if it is universal, so are all families of graph states with unbounded rank-width. In this thesis, we report progress on this problem by discussing a variety of approaches for potentially proving and disproving the universality of the comparability grid for quantum computation.

\section*{Acknowledgements}

First, I would like to thank Professor Daniel Grier for introducing me to the weird and wonderful world of quantum computing. His advice has been invaluable these past two years. He has been an incredibly kind mentor and I owe much of my growth, academically and personally, to him.

I would also like to extend this gratitude to Du (Dan) Van, the primary collaborator of this project, who has always been available to listen to my rambles. I wish him the best of luck in his future studies at NUS. I would also like to thank Jackson Morris, who has always been available to mentor and provide feedback for this project. I am also thankful for Kunal Marwaha, for helpful discussions and for sharing his notes on graph states \cite{MarwahaPrivateCorrespondence}.

I also thank my family and friends for their utmost support in all aspects of my academic journey. It is through their efforts that I may focus on my studies and pursue my aspirations without distraction.

\section*{Preface to the Second Rough Draft}

\note{Incomplete sections and comments are written in blue throughout.} 

One of the great difficulties writing the rough draft was figuring out where to stop. In other words, how much can I assume from the reader? To answer this question, the past Mathematics honors theses of Ray Tsai, Clara Woods, and Jared Hughes in particular have been most helpful; the first for its recency and treatment of graph theory, the second for its experimental and quantum focus, and the third for its approach to 3XOR and theoretical computer science. These three, among other theses, appear to take a very introductory stance, writing for a non-expert audience familiar with only the typical mathematical fields (i.e. a senior mathematics undergraduate).

Because of this, I have prepared a (perhaps excessively) extensive preliminary section and appendix. Furthermore, I may have overextended the thesis' scope, resulting in many incomplete sections in the rough draft. I was originally planning to cut out the preliminary section entirely, or perhaps place it some extended appendix. However, no prior honors thesis (to the best of my skimming) has ever used the appendix in such a manner, nor excluded so much introduction. Thus, since it is easier to cut than to add, I have left it all in.

Furthermore, the main results of the research connected to this thesis is still a work in progress --- indeed, we have not resolved the main question! Because of this, I found it challenging to create some sort of cohesive presentation of the work besides an organized diary of failed attempts. Because of this, I recognize much work needs to be done to better ``carve out'' the original research from what is known to the field. At the same time, I feel there is merit in surveying all the developments pertaining this study, even if this dwarfs the original contributions.